%% Preamble
\documentclass[paper=a4, fontsize=11pt]{scrartcl}
\usepackage[T1]{fontenc}
\usepackage{fourier}

\usepackage[english]{babel}
\usepackage[protrusion=true,expansion=true]{microtype}
\usepackage[fleqn]{amsmath}
\usepackage{amsfonts}
\usepackage{amsthm}
\usepackage[pdftex]{graphicx}
\usepackage{url}

%% Custom sectioning
\usepackage{sectsty}
\allsectionsfont{\normalfont}

%% Custom headers/footers (fancyhdr package)
\usepackage{fancyhdr}
\pagestyle{fancyplain}
\fancyhead{}                       % No page header
\fancyfoot[L]{}                    % Empty
\fancyfoot[C]{}                    % Empty
\fancyfoot[R]{\thepage}            % Pagenumbering
\renewcommand{\headrulewidth}{0pt} % Remove header underlines
\renewcommand{\footrulewidth}{0pt} % Remove footer underlines
\setlength{\headheight}{13.6pt}

%% Equation and float numbering
\numberwithin{equation}{section} % Equationnumbering: section.eq#
\numberwithin{figure}{section}   % Figurenumbering: section.fig#
\numberwithin{table}{section}    % Tablenumbering: section.tab#

%% Hyperlinks
\usepackage{hyperref}

%% Code block highlighting
\usepackage{listings}

%% Graphs
\usepackage{pgfplots}
\pgfplotsset{width=12cm,compat=1.18}

%% Maketitle metadata
\title{
    \usefont{OT1}{bch}{b}{n}
    \normalfont \normalsize \textsc{Instituto Superior Técnico} \\ [25pt]
    \huge Parallel and Distributed Computing \\ Parsim MPI Report
}
\author{
    \normalfont \normalsize
    João Gouveia - 102611, Rodrigo Arede - 102606, Gonçalo Rua - 102604 \\[-3pt] \normalsize
    \today
}
\date{}

%%% Begin document
\begin{document}
\maketitle
\section{Introduction}

For this stage of the project we explored paralellizing the particle simulator
using MPI and OpenMP. We developed several solutions, and empirically compared
them using RNL's cluster.

\section{Approach}
\subsection{Lines}

Our first approach was dividing the grid in lines, and assigning each line to a
process. This is by far the easiest approach, it is also, at least
theoretically, the worst.

\subsection{Blocks}

To minimize communication between nodes, we experimented with splitting the grid
into blocks. Each task first computes the centers of mass for the cells that
need to be communicated and sends them to the appropriate rank. While waiting to
receive the centers of mass it requires, it computes the centers of mass for the
cells that do not need to be communicated. The process for calculating particle
positions is analogous.

Block dimensions are computed dynamically. Let $ncside$ denote the size of the
grid's sides, $nprocs$ the number of tasks available to us, and $npart$ the
total number of particles. Let $size = \max(ncside / 2, 5)$ denote the height of
our blocks. Each cell has a weight corresponding to the number of particles
assigned to it, and each rank has a weight equal to the sum of the weights of
the cells assigned to it. While the rank weight is less than a threshold value,
defined as $\max(npart / nprocs, 10)$, we continue assigning vertical lines of
$size$ cells to the current rank. Once this threshold is exceeded, we advance to
the next rank in a round-robin fashion.

\subsection{Blocks And OpenMP}

After fine-tuning the parameters of our block strategy, we used OpenMP to
parallelize the loop that calculates the new particle positions. This further
reduced the running time of our program.

\section{Synchronization Concerns}

During the main loop, there is only one point where synchronization is required:
after calculating the centers of mass of all the particles, each rank must wait
to receive the centers of mass of the cells surrounding its block. Besides this,
we use OpenMP's \texttt{critical} directive to synchronize access during the
parallelized loop.

\section{Load Balancing}

Our block construction inherently balances the load by dynamically adjusting the
size of the blocks to achieve a more homogeneous distribution of particles
across ranks. Additionally, we use OpenMP's dynamic scheduling in our
parallelized loop.

\section{Performance Results}

\subsection{Methodology}

We tested our solution multiple times, using different numbers of tasks to
refine our strategies for block calculation and load balancing. The following
tests were used:

\begin{center}
    \begin{tabular}{ | c c | }
        \hline
        Test Number & Input \\
        \hline
        1 & 12672 0.05 3 10 10 \\
        \hline
        2 & 5893 0.05 3 10 10 \\
        \hline
        3 & 8555 0.05 3 10 10 \\
        \hline
        4 & 12 100 5 10000 10000 \\
        \hline
        5 & -11 3500 20 500000 10 \\
        \hline
        6 & 1 5000 100 1000000 4 \\
        \hline
        7 & 1 5000 100 1000000 100 \\
        \hline
        8 & 1 5000 20 1000000 10 \\
        \hline
        9 & 1 1000 3 10000 10000 \\
        \hline
        10 & 3 5000 50 1000000 300 \\
        \hline
        11 & 3 5000 50 1000000 500 \\ 
        \hline
        12 & -1 1000 30 100000 1000 \\
        \hline
    \end{tabular}
\end{center}

\subsection{Results}

The table below shows the execution times (in seconds) for the serial and OpenMP versions of the program:

| Test Number | 1 thread (serial) | 2 threads | 4 threads | 8 threads |
|-------------|-------------------|-----------|-----------|-----------|
| 1           | 0.0s              | 0.0s      | 0.0s      | 0.0s      |
| 2           | 0.0s              | 0.0s      | 0.0s      | 0.0s      |
| 3           | 0.0s              | 0.0s      | 0.0s      | 0.0s      |
| 4           | 53.9s             | 31.4s     | 21.2s     | 16.7s     |
| 5           | 49.7s             | 25.5s     | 13.6s     | 7.5s      |
| 6           | 1.6s              | 1.0s      | 0.8s      | 0.8s      |
| 7           | 34.9s             | 23.0s     | 18.1s     | 15.7s     |
| 8           | 57.1s             | 29.7s     | 15.6s     | 8.7s      |
| 9           | 245.9s            | 141.4s    | 89.8s     | 68.0s     |
| 10          | 289.4s            | 160.5s    | 100.1s    | 74.1s     |
| 11          | 481.9s            | 263.0s    | 171.8s    | 128.5s    |
| 12          | 81.8s             | 45.7s     | 30.4s     | 24.0s     |

\section{Conclusion}

The project successfully implemented a parallel particles simulation,
demonstrating significant speedups using OpenMP. The next step is to use MPI to
parallelize the workload across different machines, and exploring an hybrid
MPI+OpenMP approach for even better performance and scalability.

%%% End document
\end{document}
